\documentclass{umk-matematika}

\begin{document}

\maketitlepage
    {Практическая работа №11}
    {Параллельность прямых и плоскостей}
    {}

\section{Фундамент (1--4 балла)}

\textbf{Задача 1 (1 балл).} Сформулируйте признак параллельности прямой и плоскости.

\textbf{Задача 2 (2 балла).} Сформулируйте признак параллельности двух плоскостей.

\textbf{Задача 3 (3 балла).} Какие прямые в пространстве называются скрещивающимися?

\textbf{Задача 4 (4 балла).} В кубе $ABCDA_1B_1C_1D_1$ укажите рёбра, параллельные ребру $AB$.

\section{Стены (5--7 баллов)}

\textbf{Задача 5 (5 баллов).} Дана треугольная пирамида $SABC$. Точки M и N — середины рёбер SA и SB соответственно. Докажите, что прямая MN параллельна плоскости ABC.

\textbf{Задача 6 (6 баллов).} Даны две параллельные плоскости. Прямая пересекает первую плоскость. Пересечёт ли она вторую?

\textbf{Задача 7 (7 баллов).} Через концы отрезка AB и его середину M проведены параллельные прямые, пересекающие плоскость α в точках A₁, B₁, M₁. Найдите MM₁, если AA₁ = 5 см, BB₁ = 11 см, и отрезок AB не пересекает плоскость.

\section{Кровля (8--10 баллов)}

\textbf{Задача 8 (8 баллов).} В кубе $ABCDA_1B_1C_1D_1$ точка K — середина ребра AA₁, точка L — середина ребра CC₁. Докажите, что прямая KL параллельна плоскости ABCD.

\textbf{Задача 9 (9 баллов).} Даны две скрещивающиеся прямые a и b. Через точку M, не лежащую на них, проведите плоскость, параллельную обеим прямым. Всегда ли это возможно?

\textbf{Задача 10 (10 баллов).} Две балки перекрытия лежат в параллельных плоскостях. Расстояние между плоскостями — 3 м. Концы верхней балки проектируются на нижнюю плоскость. Найдите расстояние между проекциями концов верхней балки, если её длина 5 м, а угол между балками 60$^\circ$.

\newpage
\section*{Ответы}

\begin{enumerate}
  \item признак параллельности прямой и плоскости.
  \item признак параллельности плоскостей.
  \item прямые, не лежащие в одной плоскости.
  \item DC, A₁B₁, D₁C₁.
  \item $MN \parallel (ABC)$.
  \item да, пересечёт.
  \item $MM_1 = 8$ см.
  \item $KL \parallel (ABCD)$.
  \item да, всегда возможно.
  \item 5 м.
\end{enumerate}

\end{document}